\chapter{Introduction}

A well-known stylized fact about firms is how the volatility of their growth rate
depends on their size. If you measure how much a firm's growth
rate fluctuates from year to year, this volatility falls with size as $\sigma(S) \propto S^{-\beta}$, and in the data the
exponent is $\beta \approx 0.15$--$0.20$ \citep{stanley1996,
amaral1997}. Alongside this scaling, the growth-rate distribution is sharply peaked and
fat-tailed rather than Gaussian. The standard way to explain both facts is the
granular mechanism of \citet{wyartbouchaud2003} and \citet{gabaix2011}: a firm's
activity is concentrated in a few large sub-units rather than spread evenly across
many comparable ones, so a handful of dominant parts produce most of its
fluctuations, and the volatility decreases only slowly as the firm grows. Recently,
however, \citet{moran2024} confronted these models with firm data and found that they
do not hold: the size-conditioned volatility distribution has a tail far too thin to
match the granular prediction, and the growth rates stay fat-tailed even after
rescaling by each firm's own volatility. They conclude that granularity alone cannot
account for the two laws, and suggest that the missing ingredient is correlations
between a firm's parts that grow with its size. This thesis takes up that question at
the level of the whole economy: whether the interactions between firms can generate
such size-dependent correlations, and with them the empirical scaling and the heavy
tails.

The granular mechanism is best seen one level up, at the whole economy. Normally we expect the
idiosyncratic shocks hitting $N$ comparable units to average out as $1/\sqrt{N}$, so
that anything left at the macroscopic scale has to come from genuinely aggregate
shocks. \citet{gabaix2011} pointed out that this is no longer true when the units
have a fat-tailed (Zipf-like) size distribution. In that case a few big units
account for a finite fraction of the total. Their own shocks are no longer cancelled
by the others, and firm-level randomness appears directly in the aggregate. He estimated
that the largest hundred US firms alone account for about a third of the variance of
output growth. So what matters is the concentration of the size distribution, not
how many units there are. The same argument can be pushed down one level, to a
single firm seen as the sum of its own parts.

This is what the ``sub-business'' models do. A firm is treated as a collection of
internal units (divisions, products, plants), and its growth rate is the
size-weighted sum of theirs. \citet{sutton2002} called the slow size--volatility
decay a ``scaling puzzle'' and tried to explain it from the internal structure of
the firm. \citet{wyartbouchaud2003} then showed that if you give the sub-units a
heavy-tailed size distribution you get both anomalies at once: the variance of firm
growth decays too slowly with size, and the growth-rate distribution stays
non-Gaussian and fat-tailed instead of becoming Gaussian as the number of sub-units
grows. On this view the firm-level puzzle has the same granular origin as the
aggregate one: the spread in the sizes of the constituents, not their number, is what
decides how well diversification works.

\citet{moran2024} recently revisited this granular picture and tested it directly.
They first unify the models of \citet{gabaix2011} and \citet{wyartbouchaud2003} into
a single framework, writing the growth variance of a firm as
$\sigma^2 = \sigma_0^2\, H$, where $H$ is the Herfindahl index measuring how
concentrated the firm's sales are across its sub-units. A heavy-tailed sub-unit size
distribution (exponent $1 < \mu < 2$) makes $H$ decay slowly with the number of
sub-units, which would give $\beta < 1/2$ and a fat-tailed growth-rate distribution.
The framework also makes a sharper, testable prediction: once the volatility is
rescaled by its size-dependent average, its distribution across firms should be
universal, and the growth rates, rescaled by each firm's own volatility, should
become Gaussian. Neither prediction holds in the data. The rescaled volatility
distribution is indeed universal, but its tail is far too thin to match the granular
prediction, and the rescaled growth rates stay fat-tailed at every size. They
therefore conclude that granularity alone cannot explain the two laws, and argue that
the missing ingredient is correlations between a firm's parts that grow with its
size, coming from shared customers and suppliers or from firm-wide reputation shocks.

The correlations that this picture still needs sit inside the individual firm, among
a fixed, externally chosen list of sub-units, even though the sources behind them,
shared customers and suppliers, are really links between firms. In this thesis we
drop the sub-unit picture entirely. Instead of breaking a firm into internal pieces and blaming its
volatility on how concentrated they are, we model the economy directly as a network
of interacting firms: every node of the graph is a company, and the edges are the
competitive and cooperative couplings between them. The question we ask is whether
the same empirical facts can come out of this inter-firm picture, which seems to us
more economically natural, where the correlations are not imposed externally but are
produced dynamically by the interactions between companies. The standard model for
this kind of interaction is the Generalized Lotka--Volterra (GLV) system, borrowed
from theoretical ecology: each node grows logistically and is coupled to its
neighbors through a signed interaction matrix encoding competition and cooperation.
So instead of assuming a heavy-tailed firm composition, we let a network of
inter-firm couplings generate the heterogeneity on its own. The difficulty is that a
deterministic GLV system does not usually keep fluctuating: away from equilibrium,
its restoring forces just damp everything out. The interesting case is the critical
point $\mu = \mu_c$, where the interior fixed point loses stability, the restoring
forces go to zero, and you get long-lived, scale-spanning transient fluctuations.
A natural first attempt is to study this critical point: can a static GLV network
tuned to it reproduce the empirical scaling and the heavy tails on its own, as an
endogenous alternative to the granular mechanism?

Sitting right at $\mu_c$ is where a naive simulation fails. As
$\mu \to \mu_c$ the total abundance grows without bound, and for $\mu > \mu_c$ it
diverges in finite time, which forces any explicit ODE solver into tiny step sizes or
floating-point overflow. We therefore first build a numerical method that sidesteps
this rather than attacking it directly. We split the dynamics into a bounded ``shape''
part, living on the probability simplex, and an unbounded ``scale'' part carrying the
total abundance, and rescale time so that the finite-time blow-up in physical time
becomes a smooth growth in a new time variable. The shape can then be integrated
stably straight through the critical point, which also locates $\mu_c$ cheaply. With
this tool we test the static critical model across regular, exponential, and power-law
topologies. It reproduces the heavy tails of the growth-rate distribution on the
heterogeneous graphs, but it fails on the other fact: there the size--variance
relation is not a stationary law but a relaxation transient, a curve traced out as an
arbitrary initial condition relaxes onto the critical attractor, with no intrinsic
exponent to compare against the data. The static critical point thus supplies the
tails but not the slow size--variance decay, and the reason is structural: once the
shape reaches the attractor the dynamics stop, so a stationary size--variance law needs
a state that keeps fluctuating instead of relaxing.

That observation motivates the second and central part of the thesis, where we propose
a dynamical model built to fluctuate indefinitely while the economy grows. We keep the
same disordered network but let the interactions act on the average firm rather than on
absolute abundances. This single change makes the dynamics scale-invariant: the total
abundance grows exponentially without ever diverging, and in the strongly disordered
regime the composition never settles but keeps churning chaotically. In this
persistently fluctuating state the two empirical facts appear together as genuine
stationary properties, a symmetric fat-tailed growth distribution and a size--variance
relation that falls with the empirical sign, and they survive across graph
realizations, system sizes, and sampling intervals. Switching off the disorder removes
both, which places their origin in the interactions rather than in any property of a
single firm. The contribution of this thesis is therefore not to confirm that
size-dependent correlations are the missing ingredient, but to give a first-principles
dynamical model of interacting firms from which the firm-growth anomalies, and a
growing economy with persistent turnover, emerge on their own.

The rest of the thesis is organized as follows. Chapter~\ref{chap:rescaling} develops
the coordinate change and time rescaling that make the critical regime reachable, and
uses them to locate $\mu_c$ on finite graphs and against the mean-field prediction.
Chapter~\ref{chap:results} presents the static critical results across topologies,
showing that it recovers the heavy tails but that its size--variance relation is only a
relaxation transient. Chapter~\ref{chap:growing} introduces the relative model, derives
its scale-invariant growth, and measures the two firm-growth facts in its fluctuating
state, together with their robustness across system size, disorder, graph realizations,
and sampling. Chapter~\ref{chap:conclusion} draws the results together and sketches the
extensions, in particular measuring the emergent size-dependent correlations directly,
that would be needed to connect the model to the data of \citet{moran2024} and to
realistic macroeconomic dynamics.
